\documentclass[11pt]{amsart}

\usepackage[T1]{fontenc}
\usepackage[utf8]{inputenc}
\usepackage{lmodern}
\usepackage{amsmath,amssymb,amsfonts,amsthm,mathtools}
\usepackage{mathrsfs}
\usepackage{geometry}
\usepackage{booktabs}
\usepackage{microtype}
\geometry{margin=1in}
\usepackage{hyperref}
\hypersetup{
 colorlinks=true,
 linkcolor=blue,
 citecolor=blue,
 urlcolor=blue,
 pdftitle={Reciprocal Sine Products and Centered Euclidean Kernels},
 pdfauthor={Victor Casana-Giner}
}
\urlstyle{same}

\numberwithin{equation}{section}
\setcounter{tocdepth}{2}
\setcounter{secnumdepth}{2}

\newtheorem{theorem}{Theorem}[section]
\newtheorem{proposition}[theorem]{Proposition}
\newtheorem{corollary}[theorem]{Corollary}
\newtheorem{lemma}[theorem]{Lemma}
\theoremstyle{definition}
\newtheorem{definition}[theorem]{Definition}
\newtheorem{experiment}[theorem]{Numerical experiment}
\theoremstyle{remark}
\newtheorem{remark}[theorem]{Remark}

\newcommand{\C}{\mathbb C}
\newcommand{\R}{\mathbb R}
\newcommand{\eps}{\varepsilon}
\newcommand{\dd}{\,\mathrm d}
\newcommand{\Li}{\operatorname{Li}}

\title{Reciprocal Sine Products and Centered Euclidean Kernels}
\author{V\'ictor Casa\~na-Giner}
\address{Bio Logic Crop Science S.L.}
\email{\href{mailto:vcasana@biologiccropscience.com}
{vcasana@biologiccropscience.com}}
\thanks{ORCID: \href{https://orcid.org/0000-0002-8252-2608}
{0000-0002-8252-2608}.}
\date{Version 8 --- August 2026}

\begin{document}

\begin{abstract}
For integers \(n\geq3\), consider
\[
 P(n)=\prod_{k=2}^{n-1}\sin\!\left(\frac{\pi n}{k}\right).
\]
We record exact residue, gamma, Fourier, divisor-indexed, and Euclidean
forms of this varying-denominator product.  The elementary block argument
for the fractional parts \(\{n/k\}\) gives the limiting density
\(\rho(t)=\psi_1(1+t)\) and the scale
\(\log|P(n)|=-Cn+O(\sqrt n\log n)\) on the nonzero locus.  Specializing
a theorem of Saffari--Vaughan to the diagonal range improves the uniform
counting discrepancy from \(O(\sqrt n)\) to
\(O(n^{1/3}\log n)\); after the logarithmic endpoint singularities are
retained explicitly this gives, for primes,
\[
 \log|P(p)|=-Cp+O\!\left(p^{1/3}(\log p)^2\right).
\]
The same input sharpens the zero-omitted product and its primitive divisor
component, and moves the half-plane of absolute convergence of the
associated Dirichlet series from \(\Re s>3/2\) to \(\Re s>4/3\).

For prime arguments the reciprocal Fourier modes admit an exact centering,
and their partial sums are a single Dirichlet-kernel quadrature defect.
The exact Parseval identity for the rotated logarithmic product now yields
\(\|\mathscr R_p\|_2\ll p^{1/3}\log p\).  More specifically, the
unrotated remainder sees only the folded discrepancy
\(H_p(t)=D_p(t)-D_p(1-t)\).  We prove the exact formula
\[
 \log|P(p)|+Cp-2\log2
 =-\pi\int_0^{1/2}H_p(t)\cot(\pi t)\,dt,
\]
show that the extreme singular range \(t\le p^{-2/3}\) contributes only
\(O(p^{1/3}\log p)\), using the current Li--Yang divisor bound, and isolate
the remaining logarithmic loss to the mesoscopic interval
\(p^{-2/3}\le t\le1/2\).  An exact regular-grid representation of
\(H_p\) as a symmetric difference of reciprocal sawtooth sums is also
proved, with all endpoint corrections displayed.

Finally, writing \(s_k=k\|p/k\|\) separates the singular logarithm into a
smooth bounded-variation term and an integer core \(S_p=\prod s_k\).  The
core has exact shifted-divisor and gamma factorizations; the corrected
gamma formula is verified numerically and a uniform Stirling expansion
shows that the complementary-divisor sector \(q\le p^{1/3}\) can be
smoothed at total cost \(O(p^{1/3}\log p)\).  Extensive numerical audits
are included, in particular exhaustive one-sided jump maxima for the
folded discrepancy.  The observed bound \(\max|H_p|\asymp p^{1/3}\) is
stated only as a numerical conjecture.  The finite encodings of
\(\pi,\mu,\Lambda\), and \(\Psi\) remain exact reorganizations of classical
arithmetic data and do not sharpen the von Koch criterion for the Riemann
hypothesis.
\end{abstract}

\maketitle
\tableofcontents

\section{Introduction}
\label{sec:introduction}

This paper studies the diagonal, varying-denominator product
\begin{equation}\label{eq:P-def-intro}
 P(n)=\prod_{k=2}^{n-1}\sin\!\left(\frac{\pi n}{k}\right).
\end{equation}
It differs from the classical complete product with a fixed denominator
and from Sudler products, in which the frequency is fixed while the number
of factors grows.  Fixed-denominator products, partial rational sine
products, and Sudler products have extensive theories; see
\cite{Washington1997,OSullivan2016,GrepstadEtAl2019}.  The distribution of
fractional parts of \(x/n\) is likewise classical
\cite{SaffariVaughan1977}.  For the specific diagonal family
\(\{n/k\}\), the same explicit pushforward density used below was also
written out in a public 2023 derivation \cite{MikeTeX2023}.

The aim here is structural and quantitative.  We derive exact identities
for \eqref{eq:P-def-intro}, first by an elementary block argument and then
by importing the sharper diagonal discrepancy estimate of Saffari and
Vaughan.  The latter improves the proved logarithmic remainder to
\(O(p^{1/3}(\log p)^2)\) on prime arguments.  The product-specific step is
then to exploit the reflection symmetry of the logarithmic sine kernel:
the unrotated remainder depends only on a folded discrepancy, and its
extreme singular part can be treated through symmetric shifted-divisor
counts.  We also isolate an integer core and a complementary-divisor gamma
factorization, retaining all floor and residue terms until their total cost
is proved.  Separately, the Euclidean Lambert kernel, the floor-sensitive
Clausen remainder, and the exact finite encodings of classical arithmetic
functions are retained from the earlier development.  Whenever a
calculation reaches an established theorem or identity, that input is
stated with a reference.  In particular, the appearances of \(\mu\),
\(\Lambda\), \(\zeta\), and \(\Psi\) below do not furnish new information
about their size or analytic continuation.

\section{Exact forms}
\label{sec:exact-forms}

For \(n\geq3\), write
\begin{equation}\label{eq:quotient-residue}
 m_k=\left\lfloor\frac nk\right\rfloor,
 \qquad r_k=n\bmod k,
 \qquad 0\leq r_k<k.
\end{equation}

\subsection{Residues and sign}

The identity
\begin{equation}\label{eq:sine-residue}
 \sin\!\left(\frac{\pi n}{k}\right)
 =(-1)^{m_k}\sin\!\left(\frac{\pi r_k}{k}\right)
\end{equation}
separates the sign from the nonnegative residue factor.

\begin{proposition}\label{prop:residue-sign}
For every integer \(n\geq3\),
\begin{equation}\label{eq:residue-form}
 P(n)=(-1)^{n+1+\lfloor\sqrt n\rfloor}
 \prod_{k=2}^{n-1}\sin\!\left(\frac{\pi(n\bmod k)}{k}\right).
\end{equation}
For odd \(n\), the sign exponent reduces modulo two to
\(\lfloor\sqrt n\rfloor\).
\end{proposition}

\begin{proof}
The sign exponent before reduction is
\[
 \sum_{k=2}^{n-1}\left\lfloor\frac nk\right\rfloor=D(n)-n-1,
 \qquad D(n)=\sum_{j\leq n}d(j).
\]
Since \(d(j)\) is odd exactly when \(j\) is a square,
\(D(n)\equiv\lfloor\sqrt n\rfloor\pmod2\).  Formula
\eqref{eq:residue-form} follows.
\end{proof}

Taking absolute values of roots of unity gives
\begin{equation}\label{eq:finite-root-form}
 P(n)=(-1)^{n+1+\lfloor\sqrt n\rfloor}2^{-(n-2)}
 \prod_{k=2}^{n-1}\left|1-e^{2\pi i n/k}\right|.
\end{equation}

\subsection{Gamma, Fourier, and shifted-divisor forms}

Euler's reflection formula gives, whenever every factor is nonzero,
\begin{equation}\label{eq:gamma-form}
 P(n)=\frac{\pi^{n-2}}
 {\displaystyle\prod_{k=2}^{n-1}\Gamma(n/k)\Gamma(1-n/k)}.
\end{equation}
The Abel-regularized Fourier series of \(\log|2\sin\pi t|\) gives
\begin{equation}\label{eq:fourier-form}
 \log|P(n)|=-(n-2)\log2
 -\lim_{u\uparrow1}\sum_{j\geq1}\frac{u^j}{j}
 \sum_{k=2}^{n-1}\cos\!\left(\frac{2\pi jn}{k}\right).
\end{equation}
The inner sums are reciprocal-phase exponential sums of the type appearing
in the Dirichlet divisor problem; see \cite{GrahamKolesnik1991}.

For a nonzero residue \(r=n\bmod k\), one has \(k\mid n-r\) and \(k>r\),
and the converse is immediate.  Hence
\begin{equation}\label{eq:divisor-indexed}
 |P(n)|=\prod_{r=1}^{n-2}
 \prod_{\substack{k\mid n-r\\k>r}}
 \sin\!\left(\frac{\pi r}{k}\right)
\end{equation}
whenever \(P(n)\ne0\).  More locally,
\begin{equation}\label{eq:shifted-divisors}
 r_k=s\iff k\mid n-s,
 \qquad k-r_k=s\iff k\mid n+s.
\end{equation}

\section{Block distribution and exponential scale}
\label{sec:block-scale}

Partition the indices by
\(m=\lfloor n/k\rfloor\) and \(t=\{n/k\}\).  The block \(m\)
corresponds to \(n/(m+1)<k\leq n/m\).
Here \(\psi_0\) and \(\psi_1\) denote the di- and trigamma functions.
The distribution of fractional parts of \(x/n\) belongs to the classical
framework of Saffari and Vaughan \cite{SaffariVaughan1977}.  For the exact
family \(\{n/k\}\), the interval decomposition and the resulting density
\(\sum_{m\geq1}(m+t)^{-2}\) were also displayed explicitly in the 2023
public derivation \cite{MikeTeX2023}.  The following short proof records the
uniform estimate needed here.

\begin{theorem}\label{thm:block-count}
Uniformly for \(0\leq t\leq1\),
\begin{equation}\label{eq:block-count}
 \#\left\{2\leq k\leq n-1:\left\{\frac nk\right\}\leq t\right\}
 =n\bigl(\psi_0(1+t)-\psi_0(1)\bigr)+O(\sqrt n).
\end{equation}
Consequently, for every continuous \(f:[0,1]\to\C\),
\begin{equation}\label{eq:weak-density}
 \frac1n\sum_{k=2}^{n-1}f\!\left(\left\{\frac nk\right\}\right)
 \longrightarrow\int_0^1 f(t)\rho(t)\dd t,
 \qquad
 \rho(t)=\psi_1(1+t)=\sum_{m\geq1}\frac1{(m+t)^2}.
\end{equation}
The density has total mass one.
\end{theorem}

\begin{proof}
Within block \(m\), the additional condition \(\{n/k\}\leq t\) selects
an interval containing
\[
 \frac{n}{m}-\frac{n}{m+t}+O(1)
 =\frac{nt}{m(m+t)}+O(1)
\]
integers.  Summing for \(m\leq M\), and bounding both the omitted indices
and the main-term tail by \(O(n/M)\), gives total error
\(O(M+n/M)\).  Choose \(M=\lfloor\sqrt n\rfloor\) and use
\[
 \sum_{m\geq1}\frac{t}{m(m+t)}
 =\sum_{m\geq1}\left(\frac1m-\frac1{m+t}\right)
 =\psi_0(1+t)-\psi_0(1).
\]
Weak convergence gives \eqref{eq:weak-density}, and differentiation gives
the density.  Its mass is
\(\sum_{m\geq1}(1/m-1/(m+1))=1\).
\end{proof}

\medskip
\noindent\textbf{Collapse principle.}
For every continuous \(F:[0,1]\to\mathbb C\),
\begin{equation}\label{eq:collapse}
 \frac1n\sum_{k=2}^{n-1}F\!\left(\left\{\frac nk\right\}\right)
 \longrightarrow
 \int_0^1 F(t)\,\psi_1(1+t)\dd t.
\end{equation}
This is exactly the weak-convergence statement \eqref{eq:weak-density}:
writing \(u=k/n\), the grid points become Lebesgue distributed on
\((0,1]\) and \(\{n/k\}=\{1/u\}\).  Hence the diagonal geometry fixes a
single limiting pushforward class.  The observable \(F\) changes the tested
integral, but not the underlying measure.

The principle explains why several apparently different constructions in
this paper return to the same object.  The residue, gamma, Fourier, and
shifted-divisor formulas of Section~\ref{sec:exact-forms} are exact
representations of the same diagonal observable; for
\(F(t)=\log|\sin\pi t|\), the endpoint singularities require the separate
truncation estimate in Theorem~\ref{thm:scale}, but the limiting measure is
unchanged.  At the other end of the paper, the proof of
Proposition~\ref{prop:Clausen-main} again uses
\(g(u)=\operatorname{Cl}_2(2\pi\{1/u\})\), now with totient weights.
Accordingly, changing the kernel changes the tested observable and its
arithmetic weights, not the underlying fractional-part pushforward class.

The measure in \eqref{eq:weak-density} is the one-step pushforward of
Lebesgue measure by the Gauss map \(u\mapsto\{1/u\}\).  Its block masses
are \(1/(m(m+1))\), not the invariant Gauss--Kuzmin digit law; compare
\cite{Mayer1989}.  For example, the limiting proportion of negative sine
factors is
\begin{equation}\label{eq:negative-proportion}
 \sum_{\substack{m\geq1\\m\ {\rm odd}}}\frac1{m(m+1)}=\log2.
\end{equation}

The logarithmic singularities at the endpoints require a separate
estimate.

\begin{theorem}\label{thm:scale}
Along integers for which \(P(n)\ne0\),
\begin{equation}\label{eq:scale}
 \log|P(n)|=-Cn+O(\sqrt n\log n),
\end{equation}
where
\begin{align}
 C
 &=-\int_0^1\log(\sin\pi t)\,\psi_1(1+t)\dd t\label{eq:C-density}\\
 &=-\int_1^\infty\frac{\log|\sin\pi u|}{u^2}\dd u
 =-\int_0^1\log\left|\sin\frac{\pi}{u}\right|\dd u
 =0.7437948618\ldots.\label{eq:C-integrals}
\end{align}
\end{theorem}

\begin{proof}
For a nonzero factor,
\(\sin(\pi r_k/k)\geq\sin(\pi/k)\geq2/k\).  Thus the indices
\(k\leq\sqrt n\) contribute \(O(\sqrt n\log n)\).

For \(k>\sqrt n\), there are fewer than \(\sqrt n\) blocks.  Within each
block, \(k\mapsto\{n/k\}\) is monotone and
\(t\mapsto\log\sin(\pi t)\) has two monotone pieces.  Comparing the sum
with its integral therefore costs \(O(\log n)\) per block after the
singular endpoints are trimmed at the nearest admissible sample.  The
aggregate error is \(O(\sqrt n\log n)\).

With \(u=n/x\), the main integral is
\[
 n\int_{1+1/(n-1)}^{\sqrt n}
 \frac{\log|\sin\pi u|}{u^2}\dd u.
\]
Extending the upper limit to infinity costs \(O(\sqrt n)\).  At the lower endpoint the omitted integral is
\[
 n\int_1^{1+1/(n-1)}\frac{\log|\sin\pi u|}{u^2}\dd u
 =-\log n+O(1),
\]
which is absorbed by the stated \(O(\sqrt n\log n)\) remainder.  This proves
\eqref{eq:scale}.  The substitutions \(u=m+t\) and \(u\mapsto1/u\) give
\eqref{eq:C-density}--\eqref{eq:C-integrals}.
\end{proof}

\begin{remark}[Endpoint contribution]\label{rem:endpoint-log}
The lower-end calculation in the proof isolates a contribution
\(+\log n+O(1)\) when the main integral is extended down to \(u=1\).
Since the proven global error is \(O(\sqrt n\log n)\), this observation is
not asserted as a second-order asymptotic term.  In particular,
Theorem~\ref{thm:scale} claims only the first-order law \(-Cn\) with the
stated remainder.
\end{remark}

\subsection{Saffari--Vaughan refinement of the diagonal discrepancy}

The preceding proof is elementary but not quantitatively optimal.  The
standard case of Saffari--Vaughan \cite[Theorem~1 and Corollary~1.1]{SaffariVaughan1977}
contains the sharper diagonal estimate needed here.

\begin{theorem}[Saffari--Vaughan specialization]\label{thm:SV-specialization}
Uniformly for \(0\le t\le1\),
\begin{equation}\label{eq:SV-specialization}
 \#\left\{2\le k\le n-1:\left\{\frac nk\right\}\le t\right\}
 =nF(t)+O\!\left(n^{1/3}\log n\right),
 \qquad F(t)=\psi_0(1+t)-\psi_0(1).
\end{equation}
The choice of strict or non-strict endpoint convention changes the count by
at most the multiplicity of the corresponding jump and does not alter the
stated estimate.
\end{theorem}

\begin{proof}
Saffari and Vaughan normalize the count by the length \(y\).  In their
Theorem~1 the error is \(O(x^{1/3}y^{-1}\log x)\).  Set \(x=y=n\) and use
their Corollary~1.1, whose limiting function is
\[
 \sum_{m\ge1}\frac{t}{m(m+t)}
 =\psi_0(1+t)-\psi_0(1)=F(t).
\]
Multiplying the normalized estimate by \(n\) gives
\eqref{eq:SV-specialization}; removing the two endpoints \(k=1,n\) costs
only \(O(1)\).
\end{proof}

\begin{theorem}[Improved prime logarithmic scale]\label{thm:prime-third-scale}
For every prime \(p\),
\begin{equation}\label{eq:prime-third-scale}
 \log|P(p)|=-Cp+O\!\left(p^{1/3}(\log p)^2\right).
\end{equation}
\end{theorem}

\begin{proof}
Put \(g(t)=\log|2\sin\pi t|\), and let
\(E_p\ll p^{1/3}\log p\) denote the uniform discrepancy supplied by
Theorem~\ref{thm:SV-specialization}.  Choose
\(\delta=p^{-2/3}\).  On \([\delta,1-\delta]\), Stieltjes integration by
parts gives an error
\[
 O\!\left(E_p\operatorname{Var}_{[\delta,1-\delta]}g\right)
 =O\!\left(p^{1/3}(\log p)^2\right),
\]
because \(g'(t)=\pi\cot(\pi t)\) and the truncated variation is
\(O(\log(1/\delta))\).

At either endpoint, Theorem~\ref{thm:SV-specialization} gives
\(O(p\delta+E_p)\) sample points.  Primality implies
\(\|p/k\|\ge1/k\ge1/p\), so every sampled logarithm is \(O(\log p)\) in
absolute value.  The corresponding endpoint integrals are
\(O(p\delta\log p)\).  These contributions are again
\(O(p^{1/3}(\log p)^2)\).  Finally,
\[
 \sum_{k=2}^{p-1}g(p/k)-p\int_0^1g(t)\rho(t)\,dt
 =\log|P(p)|+Cp-2\log2,
\]
which proves the result.
\end{proof}

The same constant has the convergent representation
\begin{equation}\label{eq:C-series}
 C=\log2+\sum_{j\geq1}\frac1j
 \left[1-2\pi j\left(\frac\pi2-\operatorname{Si}(2\pi j)\right)\right].
\end{equation}

\begin{proposition}\label{prop:C-strict}
The constant \(C\) satisfies \(C>\log2\).
\end{proposition}

\begin{proof}
Put \(g(t)=\log(2\sin\pi t)\).  Then
\(C-\log2=-\sum_{m\geq1}I_m\), where
\[
 I_m=\int_0^{1/2}g(t)
 \bigl((m+t)^{-2}+(m+1-t)^{-2}\bigr)\dd t.
\]
On \((0,1/2)\), \(g\) is strictly increasing, the factor multiplying it
is strictly decreasing, and \(\int_0^{1/2}g(t)\dd t=0\).  The reversed
Chebyshev integral inequality gives \(I_m<0\); see
\cite{HardyLittlewoodPolya1952}.  Summing proves the assertion.
\end{proof}


\section{Centered reciprocal frequencies and quantitative rotated bounds}
\label{sec:frequency-centered}

For prime arguments, the reciprocal Fourier modes can be centered against
the same density before the singular Fourier series is summed.

\subsection{Centered modes and the exact prime remainder}

For \(j\ge1\), define
\begin{equation}\label{eq:cj-def}
 c_j=\int_0^1\cos(2\pi jt)\rho(t)\,dt,
 \qquad \rho(t)=\psi_1(1+t).
\end{equation}
Periodicity and the block expansion of \(\rho\) give
\begin{equation}\label{eq:cj-integral}
 c_j=\int_1^\infty\frac{\cos(2\pi ju)}{u^2}\,du
 =1-2\pi j\left(\frac\pi2-\operatorname{Si}(2\pi j)\right),
\end{equation}
and repeated integration by parts gives
\begin{equation}\label{eq:cj-asymptotic}
 c_j=\frac{1}{2\pi^2j^2}-\frac{3}{2\pi^4j^4}+O(j^{-6}).
\end{equation}
Thus
\begin{equation}\label{eq:C-cj}
 C=\log2+\sum_{j\ge1}\frac{c_j}{j}.
\end{equation}
For a prime \(p\), put
\begin{equation}\label{eq:Sj-Deltaj}
 S_j(p)=\sum_{k=2}^{p-1}\cos\!\left(\frac{2\pi jp}{k}\right),
 \qquad \Delta_j(p)=S_j(p)-pc_j.
\end{equation}

\begin{theorem}[Exact centered Fourier remainder]\label{thm:centered-Fourier-prime}
For every prime \(p\),
\begin{equation}\label{eq:centered-Fourier-prime}
 \log|P(p)|+Cp
 =2\log2-\lim_{u\uparrow1}\sum_{j\ge1}\frac{u^j}{j}\Delta_j(p).
\end{equation}
\end{theorem}

\begin{proof}
Insert \(S_j(p)=pc_j+\Delta_j(p)\) into
\eqref{eq:fourier-form} and use \eqref{eq:C-cj}.  The Abel regularization is
unchanged from \eqref{eq:fourier-form}.
\end{proof}

\subsection{A single quadrature defect for the cumulative modes}

Set
\begin{equation}\label{eq:BJ-def}
 B_J(p)=\sum_{j=1}^J\Delta_j(p),
 \qquad K_J(t)=\sum_{j=1}^J\cos(2\pi jt).
\end{equation}

\begin{proposition}[Exact reciprocal quadrature identity]\label{prop:BJ-quadrature}
For every prime \(p\) and every integer \(J\ge1\),
\begin{equation}\label{eq:BJ-quadrature}
 B_J(p)=\sum_{k=2}^{p-1}K_J\!\left(\frac pk\right)
 -\int_0^pK_J\!\left(\frac px\right)\,dx,
\end{equation}
where the integral at \(x=0\) is improper.
\end{proposition}

\begin{proof}
By definition,
\[
 B_J(p)=\sum_{k=2}^{p-1}K_J(p/k)
 -p\int_0^1K_J(t)\rho(t)\,dt.
\]
Since \(\rho(t)=\sum_{m\ge1}(m+t)^{-2}\) and \(K_J\) is one-periodic,
\[
 \int_0^1K_J(t)\rho(t)\,dt
 =\int_1^\infty\frac{K_J(u)}{u^2}\,du.
\]
The substitution \(x=p/u\) gives \eqref{eq:BJ-quadrature}.
\end{proof}

\begin{remark}\label{rem:BJ-current-role}
The quadrature identity is exact, but a formerly considered bound
\(|B_J(p)|\ll\sqrt{pJ}\) would only imply an \(O(\sqrt p)\) pointwise
remainder.  Theorem~\ref{thm:prime-third-scale} is asymptotically stronger,
so such a square-root criterion is no longer the quantitative target of
this paper.  In particular, no complete-period average in the Fourier index
is used to infer a short-frequency estimate.
\end{remark}

\subsection{Exact energy identity and the strengthened rotated theorem}

For a prime \(p\), define
\begin{equation}\label{eq:Np-Dp}
 N_p(t)=\#\left\{2\le k\le p-1:\left\{\frac pk\right\}\le t\right\},
 \qquad D_p(t)=N_p(t)-pF(t)+2t.
\end{equation}
Then \(D_p(0)=D_p(1)=0\), and
Theorem~\ref{thm:SV-specialization} gives
\begin{equation}\label{eq:Dp-uniform}
 \|D_p\|_\infty\ll p^{1/3}\log p.
\end{equation}
For \(j\ge1\), put
\begin{equation}\label{eq:Aj-def}
 A_j(p)=\sum_{k=2}^{p-1}e^{2\pi i jp/k}
 -p\int_0^1e^{2\pi ijt}\rho(t)\,dt.
\end{equation}

\begin{lemma}[Exact energy identity]\label{lem:Aj-energy}
With \(\overline D_p=\int_0^1D_p(t)\,dt\),
\begin{equation}\label{eq:Aj-energy}
 \sum_{j\ge1}\frac{|A_j(p)|^2}{j^2}
 =2\pi^2\int_0^1|D_p(t)-\overline D_p|^2\,dt
 \ll p^{2/3}(\log p)^2.
\end{equation}
\end{lemma}

\begin{proof}
The signed measure
\[
 d\mu_p(t)=\sum_{k=2}^{p-1}\delta_{\{p/k\}}(dt)
 -p\rho(t)\,dt+2\,dt
\]
has cumulative function \(D_p\) and total mass zero.  Thus
\[
 A_j(p)=-2\pi i j\int_0^1D_p(t)e^{2\pi ijt}\,dt.
\]
Parseval applied to \(D_p-\overline D_p\) proves the identity; the final
bound follows from \eqref{eq:Dp-uniform}.
\end{proof}

For \(\theta\in\mathbb R/\mathbb Z\), define
\begin{equation}\label{eq:rotated-remainder}
 \mathscr R_p(\theta)=
 \sum_{k=2}^{p-1}\log\left|2\sin\pi\left(\frac pk+\theta\right)\right|
 -p\int_0^1\log|2\sin\pi(t+\theta)|\rho(t)\,dt.
\end{equation}

\begin{theorem}[Quantitative rotated \(L^2\) bound]\label{thm:rotated-L2}
For every prime \(p\),
\begin{equation}\label{eq:rotated-L2}
 \int_0^1|\mathscr R_p(\theta)|^2\,d\theta
 =\frac12\sum_{j\ge1}\frac{|A_j(p)|^2}{j^2}
 =\pi^2\int_0^1|D_p(t)-\overline D_p|^2\,dt
 \ll p^{2/3}(\log p)^2.
\end{equation}
Moreover,
\begin{equation}\label{eq:theta-zero-original}
 \mathscr R_p(0)=\log|P(p)|+Cp-2\log2.
\end{equation}
\end{theorem}

\begin{proof}
The Fourier coefficients of
\(g(t)=\log|2\sin\pi t|\) are
\(\widehat g(j)=-1/(2|j|)\) for \(j\ne0\).  Hence
\(\widehat{\mathscr R_p}(j)=-A_j(p)/(2j)\) for \(j\ge1\), and Parseval
with Lemma~\ref{lem:Aj-energy} proves \eqref{eq:rotated-L2}.  The value at
\(\theta=0\) follows from
\(\int_0^1g(t)\rho(t)\,dt=-C+\log2\).
\end{proof}

\begin{corollary}\label{cor:rotated-Chebyshev}
For every \(\lambda>0\),
\begin{equation}\label{eq:rotated-exceptional}
 \operatorname{meas}\left\{\theta:\ |\mathscr R_p(\theta)|>
 \lambda p^{1/3}\log p\right\}\ll\lambda^{-2}.
\end{equation}
\end{corollary}

\begin{proof}
Apply Chebyshev's inequality to \eqref{eq:rotated-L2}.
\end{proof}

\section{Folded discrepancy and localization of the logarithmic loss}
\label{sec:folded-discrepancy}

The logarithmic sine kernel is reflection invariant, so the unrotated
remainder does not see the full discrepancy \(D_p\).

\begin{definition}\label{def:Hp}
For \(0\le t\le1/2\), put
\begin{equation}\label{eq:Hp-def}
 H_p(t)=D_p(t)-D_p(1-t).
\end{equation}
Also define
\begin{equation}\label{eq:Mp-Ffold}
 M_p(t)=\#\left\{2\le k\le p-1:\left\|\frac pk\right\|\le t\right\},
 \qquad F_{\rm fold}(t)=1+F(t)-F(1-t).
\end{equation}
Away from the finitely many boundary points, and with the corresponding
one-sided convention at a boundary,
\begin{equation}\label{eq:Hp-Mp}
 H_p(t)=M_p(t)-pF_{\rm fold}(t)+4t.
\end{equation}
In particular \(H_p(0)=H_p(1/2)=0\).
\end{definition}

\begin{theorem}[Exact folded Stieltjes formula]\label{thm:folded-Stieltjes}
For every prime \(p\),
\begin{equation}\label{eq:folded-Stieltjes}
 \log|P(p)|+Cp-2\log2
 =-\pi\int_0^{1/2}H_p(t)\cot(\pi t)\,dt,
\end{equation}
where the integral is improper at zero.
\end{theorem}

\begin{proof}
With \(g(t)=\log|2\sin\pi t|\), the left side is
\(\int_0^1 g\,dD_p\).  Since \(D_p(0)=D_p(1)=0\), Stieltjes integration
by parts gives
\[
 -\pi\int_0^1D_p(t)\cot(\pi t)\,dt.
\]
On the second half substitute \(t\mapsto1-t\) and use
\(\cot(\pi(1-t))=-\cot(\pi t)\).  The two halves combine into
\eqref{eq:folded-Stieltjes}.
\end{proof}

Let
\begin{equation}\label{eq:theta-star}
 \theta_*=0.3144831759741\ldots,
\end{equation}
the exponent in the current Li--Yang bound
\(\Delta(x)=O_\eps(x^{\theta_*+\eps})\) for the Dirichlet divisor problem
\cite{LiYang2023}.

\begin{proposition}[Extreme folded range]\label{prop:extreme-folded-range}
For every fixed \(\eps>0\), uniformly for \(0<t\le1/2\),
\begin{equation}\label{eq:Mp-divisor-bound}
 M_p(t)\ll pt\log p+p^{\theta_*+\eps}.
\end{equation}
Consequently,
\begin{equation}\label{eq:Hp-divisor-bound}
 H_p(t)\ll pt\log p+p^{\theta_*+\eps},
\end{equation}
and, for sufficiently small fixed \(\eps>0\),
\begin{equation}\label{eq:extreme-integral}
 \int_0^{p^{-2/3}}|H_p(t)|\,|\cot(\pi t)|\,dt
 \ll p^{1/3}\log p.
\end{equation}
\end{proposition}

\begin{proof}
For each counted denominator let
\[
 s_k=k\left\|\frac pk\right\|=\min(p\bmod k,k-(p\bmod k)).
\]
Then \(s_k\) is a positive integer, \(s_k\le pt\), and
\(k\mid p-s_k\) or \(k\mid p+s_k\).  Hence
\[
 M_p(t)\le\sum_{1\le s\le pt}\bigl(d(p-s)+d(p+s)\bigr).
\]
Apply the Li--Yang estimate for the summatory divisor function at the two
endpoints of the interval \([p-pt,p+pt]\).  This gives
\eqref{eq:Mp-divisor-bound}.  Since \(F_{\rm fold}(0)=0\) and
\(F_{\rm fold}'(t)=\rho(t)+\rho(1-t)\) is bounded on \([0,1/2]\),
\eqref{eq:Hp-Mp} gives \eqref{eq:Hp-divisor-bound}.

For \(t<1/p\), \(M_p(t)=0\), hence \(H_p(t)=O(pt)\).  On
\([1/p,p^{-2/3}]\), use \eqref{eq:Hp-divisor-bound} and
\(|\cot(\pi t)|\ll1/t\).  The \(pt\log p\) term integrates to
\(O(p^{1/3}\log p)\), while the divisor-error term contributes
\(O(p^{\theta_*+\eps}\log p)=o(p^{1/3}\log p)\) if
\(\eps<1/3-\theta_*\).
\end{proof}

\begin{corollary}[Mesoscopic location of the remaining logarithm]\label{cor:mesoscopic-log}
Theorem~\ref{thm:prime-third-scale} can be recovered from
\eqref{eq:folded-Stieltjes}, and the only range in that argument capable of
producing the second logarithm is
\begin{equation}\label{eq:mesoscopic-range}
 p^{-2/3}\le t\le\frac12.
\end{equation}
Indeed,
\begin{equation}\label{eq:mesoscopic-current-bound}
 \int_{p^{-2/3}}^{1/2}|H_p(t)|\,|\cot(\pi t)|\,dt
 \ll p^{1/3}(\log p)^2
\end{equation}
by \(\|H_p\|_\infty\ll p^{1/3}\log p\).
\end{corollary}

\subsection{Exact regular-grid representation}

To remove endpoint ambiguity, define the strict count
\begin{equation}\label{eq:N-strict}
 N_p^{<}(t)=\#\left\{2\le k\le p-1:\left\{\frac pk\right\}<t\right\},
 \quad
 D_p^{<}(t)=N_p^{<}(t)-pF(t)+2t,
\end{equation}
and \(H_p^{<}(t)=D_p^{<}(t)-D_p^{<}(1-t)\).

\begin{theorem}[Exact reciprocal sawtooth identity]\label{thm:Hp-grid}
For \(0<t<1\),
\begin{equation}\label{eq:Hp-grid}
\begin{aligned}
 H_p^{<}(t)
={}&\sum_{m=1}^{p}
 \left(\left\{\frac{p}{m+t}\right\}
       -\left\{\frac{p}{m+1-t}\right\}\right)\\
&+p\left[\psi_0(p+2-t)-\psi_0(p+1+t)\right]+4t-2.
\end{aligned}
\end{equation}
At a jump, switching from \(<\) to \(\le\) adds exactly the corresponding
jump multiplicity.  Away from jumps, \(H_p^{<}=H_p\).
\end{theorem}

\begin{proof}
Include temporarily the two endpoint denominators \(k=1,p\).  For strict
counting, partition by \(m=\lfloor p/k\rfloor\):
\[
 \#\left\{1\le k\le p:\left\{\frac pk\right\}<t\right\}
 =\sum_{m=1}^{p}
 \left(\left\lfloor\frac pm\right\rfloor
       -\left\lfloor\frac{p}{m+t}\right\rfloor\right).
\]
Use \(\lfloor x\rfloor=x-\{x\}\) and
\[
 \sum_{m=1}^{p}\left(\frac1m-\frac1{m+t}\right)
 =\psi_0(p+1)-\psi_0(p+1+t)+F(t).
\]
After removing \(k=1,p\), this yields an exact formula for
\(D_p^{<}(t)\).  Subtract the same formula with \(t\) replaced by \(1-t\):
the terms \(\{p/m\}\) and \(\psi_0(p+1)\) cancel, leaving
\eqref{eq:Hp-grid}.
\end{proof}

The smooth correction in \eqref{eq:Hp-grid} satisfies
\begin{equation}\label{eq:Hp-grid-smooth-asymp}
 p\left[\psi_0(p+2-t)-\psi_0(p+1+t)\right]+4t-2
 =2t-1+O(p^{-1})
\end{equation}
uniformly for \(0\le t\le1\).  This explains why omitting it in numerical
reconstructions produces a bounded, \(t\)-dependent residual rather than a
divergent error.

\begin{lemma}[Where the continuous maximum occurs]\label{lem:Hp-jumps}
On every open interval between consecutive jump points of \(M_p\),
\(H_p\) is strictly decreasing.  Hence the supremum of \(|H_p|\) on
\([0,1/2]\) is attained as a one-sided value at a jump.
\end{lemma}

\begin{proof}
Between jumps, \eqref{eq:Hp-Mp} gives
\[
 H_p'(t)=4-p\bigl(\rho(t)+\rho(1-t)\bigr).
\]
The function in parentheses is convex and symmetric about \(1/2\), hence
its minimum is
\[
 2\psi_1(3/2)=\pi^2-8>4/3.
\]
Thus \(H_p'(t)<0\) for every prime \(p\ge3\).
\end{proof}

\begin{remark}[Numerical conjecture]\label{rem:Hp-conjecture}
The exhaustive jump computations in Section~\ref{sec:numerical-audit}
suggest
\begin{equation}\label{eq:Hp-conjecture}
 \sup_{p^{-2/3}\le t\le1/2}|H_p(t)|\ll p^{1/3}.
\end{equation}
No proof of \eqref{eq:Hp-conjecture} is given.  If it holds, then
Theorem~\ref{thm:folded-Stieltjes} and
Proposition~\ref{prop:extreme-folded-range} imply the sharper product bound
\begin{equation}\label{eq:Hp-conjecture-consequence}
 \log|P(p)|=-Cp+O(p^{1/3}\log p).
\end{equation}
\end{remark}

\section{Integer singular core and complementary-divisor factorization}
\label{sec:integer-core}

The logarithmic singularity can be separated from a smooth
bounded-variation term.  For a prime \(p\), set
\begin{equation}\label{eq:sk-def}
 s_k=k\left\|\frac pk\right\|
 =\min(p\bmod k,k-(p\bmod k)),
 \qquad 2\le k\le p-1,
\end{equation}
and
\begin{equation}\label{eq:h-def}
 h(u)=\log\frac{\sin\pi u}{\pi u},\qquad 0<u\le\frac12,
 \quad h(0)=0.
\end{equation}
Then \(h'(u)=\pi\cot(\pi u)-1/u<0\) and
\begin{equation}\label{eq:h-variation}
 \operatorname{Var}_{[0,1/2]}h=\log\frac\pi2.
\end{equation}

\begin{proposition}[Exact smooth/integer split]\label{prop:integer-split}
Let
\begin{equation}\label{eq:Sp-def}
 S_p=\prod_{k=2}^{p-1}s_k.
\end{equation}
Then
\begin{equation}\label{eq:integer-split}
 \log|P(p)|=(p-2)\log\pi+\log S_p-\log((p-1)!)
 +\sum_{k=2}^{p-1}h\!\left(\frac{s_k}{k}\right).
\end{equation}
Define
\begin{align}
 J_0&=\int_0^{1/2}\log u\,[\rho(u)+\rho(1-u)]\,du,
 \label{eq:J0-def}\\
 H_0&=\int_0^{1/2}h(u)\,[\rho(u)+\rho(1-u)]\,du.
 \label{eq:H0-def}
\end{align}
Then
\begin{equation}\label{eq:C-JH}
 \log\pi+J_0+H_0=-C,
\end{equation}
and, with
\begin{equation}\label{eq:Acore-def}
 \mathcal A(p)=\log S_p-\log((p-1)!)-pJ_0,
\end{equation}
we have
\begin{equation}\label{eq:Acore-isolation}
 \log|P(p)|+Cp
 =\mathcal A(p)-2\log\pi+O(p^{1/3}\log p).
\end{equation}
\end{proposition}

\begin{proof}
The identity
\[
 \log\sin(\pi u)=\log\pi+\log u+h(u)
\]
with \(u=s_k/k\) gives \eqref{eq:integer-split}.  Folding the density gives
\eqref{eq:C-JH}.  Since \(h\) has bounded variation independent of \(p\),
Theorem~\ref{thm:SV-specialization} applied to the folded distribution gives
\[
 \sum_{k=2}^{p-1}h(s_k/k)=pH_0+O(p^{1/3}\log p).
\]
Substitution proves \eqref{eq:Acore-isolation}.
\end{proof}

Numerically,
\begin{equation}\label{eq:JH-numerical}
 J_0=-1.750632969245044\ldots,
 \qquad H_0=-0.137891778443756\ldots.
\end{equation}
Thus every fluctuation larger than \(p^{1/3}\log p\) in the proved
\(O(p^{1/3}(\log p)^2)\) remainder must reside in the explicit integer
core \(\mathcal A(p)\).

\subsection{A symmetric shifted-divisor count}

For an integer \(S\) with \(1\le S<(p-1)/2\), let
\begin{equation}\label{eq:NpS-def}
 \mathcal N_p(S)=\#\{2\le k\le p-1:s_k\le S\}.
\end{equation}

\begin{proposition}\label{prop:NpS-exact}
One has the exact identity
\begin{equation}\label{eq:NpS-exact}
\begin{aligned}
 \mathcal N_p(S)
={}&S-2+\sum_{n=p-S}^{p+S}d(n)\\
&-\sum_{k=1}^{2S}
 \left(\left\lfloor\frac{p+S}{k}\right\rfloor
       -\left\lfloor\frac{p-S-1}{k}\right\rfloor\right).
\end{aligned}
\end{equation}
\end{proposition}

\begin{proof}
All \(2\le k\le2S\) satisfy \(s_k\le k/2\le S\).  For \(k>2S\), the
condition \(s_k\le S\) is equivalent to the existence of a unique integer
\(q\) such that \(|p-qk|\le S\).  Counting the pairs \((q,k)\) by the
integer \(n=qk\in[p-S,p+S]\), and then subtracting the divisors
\(k\le2S\), gives \eqref{eq:NpS-exact}; the divisor \(k=1\) must be
included in the correction sum.
\end{proof}

\subsection{Complementary divisors and the corrected gamma factorization}

Every \(k\ge3\) belongs to exactly one of the branches
\(p-s_k=qk\) or \(p+s_k=qk\).  The condition \(k>2s_k\) gives
\begin{equation}\label{eq:Sp-progression-log}
\begin{aligned}
 \log S_p
={}&\sum_{q\ge1}\ \sum_{\substack{s\equiv p\ (q)\\(2q+1)s<p}}\log s\\
&+\sum_{q\ge2}\ \sum_{\substack{s\equiv-p\ (q)\\(2q-1)s<p}}\log s.
\end{aligned}
\end{equation}
The exceptional denominator \(k=2\) has \(s_2=1\) and does not affect the
logarithm.

For \(q\ge2\), put
\begin{equation}\label{eq:gamma-parameters}
 a_q=p\bmod q,\quad \alpha_q=\frac{a_q}{q},\quad
 X_q^- =\left\lfloor\frac{p-1}{2q+1}\right\rfloor,\quad
 X_q^+ =\left\lfloor\frac{p-1}{2q-1}\right\rfloor,
\end{equation}
and
\begin{align}
 n_-(q)&=\max\left(0,1+\left\lfloor\frac{X_q^--a_q}{q}\right\rfloor\right),
 \label{eq:nminus}\\
 n_+(q)&=\max\left(0,1+\left\lfloor\frac{X_q^+-(q-a_q)}{q}\right\rfloor\right).
 \label{eq:nplus}
\end{align}

\begin{proposition}[Corrected gamma factorization]\label{prop:gamma-core}
For \(q\ge2\), define
\begin{equation}\label{eq:Gq-correct}
\begin{aligned}
 \mathcal G_q(p)
={}&q^{n_-(q)+n_+(q)}\frac{|\sin(\pi p/q)|}{\pi}\\
&\times
 \Gamma\!\left(n_-(q)+\alpha_q\right)
 \Gamma\!\left(n_+(q)+1-\alpha_q\right).
\end{aligned}
\end{equation}
Then
\begin{equation}\label{eq:Sp-gamma}
 S_p=\left\lfloor\frac{p-1}{3}\right\rfloor!
 \prod_{q=2}^{(p-1)/2}\mathcal G_q(p).
\end{equation}
\end{proposition}

\begin{proof}
The first progression in \eqref{eq:Sp-progression-log} is
\(a_q,a_q+q,\ldots,a_q+(n_--1)q\), whose product is
\[
 q^{n_-}\frac{\Gamma(n_-+\alpha_q)}{\Gamma(\alpha_q)}.
\]
The second progression starts at \(q-a_q\) and contributes
\[
 q^{n_+}\frac{\Gamma(n_++1-\alpha_q)}{\Gamma(1-\alpha_q)}.
\]
Euler reflection gives
\[
 \frac1{\Gamma(\alpha_q)\Gamma(1-\alpha_q)}
 =\frac{\sin(\pi\alpha_q)}{\pi}
 =\frac{|\sin(\pi p/q)|}{\pi}.
\]
For \(q=1\), the first branch contributes
\(1\cdot2\cdots\lfloor(p-1)/3\rfloor\).  Multiplication gives
\eqref{eq:Sp-gamma}.
\end{proof}

\begin{proposition}[Uniform low-complementary-divisor smoothing]\label{prop:low-q-stirling}
Let \(2\le Q\le p^{1/3}\).  Put
\[
 Y_q^- =\frac{p-1}{q(2q+1)},\qquad
 Y_q^+ =\frac{p-1}{q(2q-1)}
\]
and
\begin{equation}\label{eq:Mq-smooth}
\begin{aligned}
 \mathcal M_q(p)
={}&\left(Y_q^--\frac12\right)
 \log\frac{p-1}{2q+1}-Y_q^-\\
&+\left(Y_q^+-\frac12\right)
 \log\frac{p-1}{2q-1}-Y_q^+.
\end{aligned}
\end{equation}
Then
\begin{equation}\label{eq:low-q-stirling}
 \sum_{q=2}^{Q}\log\mathcal G_q(p)
 =\sum_{q=2}^{Q}\mathcal M_q(p)
 +O\!\left(Q\log p+\frac{Q^3}{p}\right).
\end{equation}
In particular, at \(Q=\lfloor p^{1/3}\rfloor\),
\begin{equation}\label{eq:low-q-third}
 \sum_{q\le p^{1/3}}\log\mathcal G_q(p)
 =\sum_{q\le p^{1/3}}\mathcal M_q(p)+O(p^{1/3}\log p).
\end{equation}
\end{proposition}

\begin{proof}
For \(q\le Q\), the two large gamma arguments in
\eqref{eq:Gq-correct} have the form
\[
 z_q^-=Y_q^-+\eta_q^-,\qquad
 z_q^+=Y_q^++\eta_q^+,
 \qquad 0<\eta_q^\pm\le1,
\]
with \(Y_q^\pm\asymp p/q^2\).  Stirling's formula, uniformly in this
range, gives
\[
 \log\Gamma(z)=\left(z-\frac12\right)\log z-z
 +\frac12\log(2\pi)+O(z^{-1}).
\]
Moving \(z_q^\pm\) to \(Y_q^\pm\) costs \(O(\log p)\) per \(q\), while
the two Stirling remainders sum to
\(O(\sum_{q\le Q}q^2/p)=O(Q^3/p)\).  Finally
\(|\sin(\pi p/q)|\gg1/q\), so the explicit sine factor costs only
\(O(\log q)\) per \(q\), already contained in \(O(Q\log p)\).
\end{proof}

\begin{remark}\label{rem:gamma-limitation}
The gamma factorization is exact and the low-\(q\) smoothing is a genuine
localization statement.  It does not, by itself, improve
Theorem~\ref{thm:prime-third-scale}: substituting the reflected sine
factors back into the complete product reconstructs a complementary part of
the original product.  Any sharper global estimate must therefore control
the remaining high-\(q\) shifted-divisor sector or the folded discrepancy,
rather than rely on gamma reflection alone.
\end{remark}

\section{Primitive divisor decomposition}
\label{sec:primitive}

To obtain finite logarithms for every integer, omit precisely the zero
divisor factors:
\begin{equation}\label{eq:Ptilde}
 \widetilde P(n)=
 \prod_{\substack{2\leq k\leq n-1\\k\nmid n}}
 \left|\sin\!\left(\frac{\pi n}{k}\right)\right|,
 \qquad
 \widetilde\lambda(n)=\log\widetilde P(n).
\end{equation}
Define \(L(1)=0\) and, for \(d\geq2\),
\begin{equation}\label{eq:L-def}
 L(d)=\sum_{\substack{2\leq k\leq d-1\\(k,d)=1}}
 \log\sin\!\left(\frac{\pi(d\bmod k)}{k}\right).
\end{equation}

\begin{theorem}\label{thm:primitive}
For every \(n\geq1\),
\begin{equation}\label{eq:primitive}
 \widetilde\lambda(n)=\sum_{d\mid n}L(d),
 \qquad L=\mu*\widetilde\lambda.
\end{equation}
\end{theorem}

\begin{proof}
For \(k\nmid n\), put \(g=(n,k)\), \(d=n/g\), and \(k=gk'\).
Then \((d,k')=1\), \(2\leq k'<d\), and
\(\{n/k\}=\{d/k'\}\).  Conversely, a divisor \(d\mid n\) and a
primitive index \(2\leq k'<d\) reconstruct \(k=(n/d)k'\).  This is a
bijection between the factors on the two sides.  The second formula is the
standard M\"obius inversion formula \cite[Chapter~2]{Apostol1976}.
\end{proof}

The same bijection works for every function of \(\{n/k\}\); thus the
divisor convolution itself is not sine-specific.  With
\begin{equation}\label{eq:RL-def}
 R_L(d)=L(d)+C\varphi(d),
 \qquad \widetilde R(n)=\widetilde\lambda(n)+Cn,
\end{equation}
the classical identity \(\sum_{d\mid n}\varphi(d)=n\) gives
\begin{equation}\label{eq:R-convolution}
 \widetilde R=\mathbf1*R_L.
\end{equation}

\begin{corollary}\label{cor:tilde-scale}
Uniformly for \(n\geq2\),
\begin{equation}\label{eq:tilde-scale}
 \widetilde\lambda(n)=-Cn+O\!\left(n^{1/3}(\log n)^2\right),
 \qquad R_L(n)=O_\eps(n^{1/3+\eps}).
\end{equation}
\end{corollary}

\begin{proof}
Apply Theorem~\ref{thm:SV-specialization} to the full reciprocal
fractional-part multiset and delete the zero factors \(k\mid n\).  This
changes the distribution function by at most \(d(n)\), and every remaining
logarithm has size \(O(\log n)\).  Thus the proof of
Theorem~\ref{thm:prime-third-scale} gives
\(O(n^{1/3}(\log n)^2+d(n)\log n)\), and the divisor term is absorbed.
M\"obius inversion in \eqref{eq:R-convolution} gives
\[
 |R_L(n)|\le\sum_{d\mid n}|\widetilde R(n/d)|
 \ll n^{1/3}(\log n)^2\sum_{d\mid n}d^{-1/3}
 =O_\eps(n^{1/3+\eps}).
\]
\end{proof}

For \(\Re s>4/3\), absolute convergence and
\eqref{eq:R-convolution} give the standard Dirichlet-series identity
\begin{equation}\label{eq:AB}
 A(s)=\zeta(s)B(s),
 \qquad
 A(s)=\sum_{n\geq1}\frac{\widetilde R(n)}{n^s},
 \quad B(s)=\sum_{n\geq1}\frac{R_L(n)}{n^s}.
\end{equation}
This is exactly the Dirichlet-series image of a divisor convolution; no
continuation of \(A\) or \(B\) is asserted; compare
\cite[Chapter~2]{Apostol1976}.

\section{Euclidean Lambert expansion and centered kernel}
\label{sec:centered-kernel}

Put
\begin{equation}\label{eq:ell-H}
 \ell(x)=\log\sin(\pi x),
 \qquad
 H_k(z)=\sum_{\substack{1\leq r<k\\(r,k)=1}}
 \ell\!\left(\frac rk\right)z^r.
\end{equation}

\begin{theorem}\label{thm:Lambert}
For \(|z|<1\),
\begin{equation}\label{eq:Lambert}
 \mathcal G(z):=\sum_{n\geq2}L(n)z^n
 =\sum_{k\geq2}\frac{z^k}{1-z^k}H_k(z).
\end{equation}
Moreover, \(H_k(z)=z^kH_k(z^{-1})\) as an identity of Laurent
polynomials.
\end{theorem}

\begin{proof}
Write \(n=mk+r\) in every summand of \(L(n)\), with \(m\geq1\),
\(1\leq r<k\), and \((r,k)=1\).  Absolute convergence for \(|z|<1\)
allows the rearrangement
\[
 \sum_{n\geq2}L(n)z^n
 =\sum_{k\geq2}\sum_{\substack{1\leq r<k\\(r,k)=1}}
 \ell(r/k)z^r\sum_{m\geq1}z^{mk}.
\]
The involution \(r\mapsto k-r\) proves the reflection identity.
\end{proof}

\begin{theorem}\label{thm:centered}
For \(u\in\R\), define
\begin{equation}\label{eq:Vk}
 V_k(u)=e^{iku/2}H_k(e^{-iu})-H_k(1).
\end{equation}
Then \(V_k\) is real, even, and nonnegative, with
\begin{equation}\label{eq:Vk-positive}
 V_k(u)=2\sum_{\substack{1\leq r<k/2\\(r,k)=1}}
 \left[-\ell\!\left(\frac rk\right)\right]
 \left[1-\cos\!\left(\left(\frac k2-r\right)u\right)\right].
\end{equation}
The kernel \(V_k(u-v)\) is conditionally negative definite.  Thus, if
\(\sum_jc_j=0\),
\begin{equation}\label{eq:CND}
 \sum_{i,j}c_i\overline{c_j}V_k(u_i-u_j)\leq0.
\end{equation}
On the Laplace axis,
\begin{equation}\label{eq:Delta}
 \Delta_k(t)=e^{kt/2}H_k(e^{-t})-H_k(1)\leq0
 \qquad(t\geq0).
\end{equation}
Moreover, for every \(j\geq1\),
\begin{equation}\label{eq:Delta-even-derivatives}
 \Delta_k^{(2j)}(0)
 =2\sum_{\substack{1\leq r<k/2\\(r,k)=1}}
 \ell(r/k)\left(\frac k2-r\right)^{2j}\leq0,
\end{equation}
with strict inequality whenever the reduced-residue set contains a point
distinct from \(k/2\).
\end{theorem}

\begin{proof}
Pair \(r\) with \(k-r\).  Since \(\ell(r/k)<0\) away from the central
point, the pairing gives \eqref{eq:Vk-positive}.  Writing
\(a_{k,r}=k/2-r\) and \(w_{k,r}=-2\ell(r/k)>0\), the quadratic form in
\eqref{eq:CND} equals
\[
 -\frac12\sum_r w_{k,r}
 \left(\left|\sum_jc_je^{ia_{k,r}u_j}\right|^2
 +\left|\sum_jc_je^{-ia_{k,r}u_j}\right|^2\right),
\]
because the constant term is zero.  This is the standard negative-type
form of Schoenberg \cite{Schoenberg1938}.  Replacing \(u\) by \(-it\)
gives
\[
 \Delta_k(t)=2\sum_r\ell(r/k)
 \bigl(\cosh(a_{k,r}t)-1\bigr)\leq0.
\]
Differentiating this finite sum at the origin gives
\eqref{eq:Delta-even-derivatives}.
\end{proof}

\begin{remark}[Universality of the centered sign]\label{rem:universal-centered}
The proof of Theorem~\ref{thm:centered} uses only reflection symmetry
under \(r\mapsto k-r\) and the inequality \(\ell(r/k)<0\) away from the
central point.  Consequently, \eqref{eq:Vk-positive}, conditional negative
definiteness, \eqref{eq:Delta}, and \eqref{eq:Delta-even-derivatives}
remain valid verbatim if the reduced residues are replaced by any subset
\(S_k\subseteq\{1,\ldots,k-1\}\) satisfying
\(r\in S_k\Longleftrightarrow k-r\in S_k\).  Coprimality therefore plays
no role in these signs.

For the reduced-residue block, the arithmetic content is instead carried
by the positive spectral measure
\begin{equation}\label{eq:nu-k}
 \nu_k=\sum_{\substack{1\leq r<k/2\\(r,k)=1}}
 \left[-\ell\!\left(\frac rk\right)\right]
 \left(\delta_{k/2-r}+\delta_{-(k/2-r)}\right),
\end{equation}
for which
\[
 V_k(u)=\int_{\mathbb R}\bigl(1-\cos(\xi u)\bigr)\,d\nu_k(\xi).
\]
Thus the atoms \(k/2-r\) and their weights, not positivity by itself,
retain the reduced-residue arithmetic.  In particular, the signs
\(\mathcal D(t)<0\) and \(\kappa>0\) below are termwise signs; the
asymptotic size and spectral distribution are the nontrivial information.
\end{remark}

The complete value of a block is classical:
\begin{equation}\label{eq:H-classical}
 H_k(1)=\Lambda(k)-\varphi(k)\log2.
\end{equation}
Indeed, this is \(\Phi_k(1)=e^{\Lambda(k)}\), discussed further in
Section~\ref{sec:recoveries}.  The centered kernel subtracts this total
mass before retaining the frequencies \(k/2-r\).

Define
\begin{equation}\label{eq:mathcalD}
 \mathcal D(t)=\sum_{k\geq2}
 \frac{e^{-3kt/2}}{1-e^{-kt}}\Delta_k(t),
 \qquad t>0.
\end{equation}

\begin{proposition}\label{prop:D-scale}
As \(t\downarrow0\),
\begin{equation}\label{eq:D-scale}
 \mathcal D(t)\sim-\frac{\kappa}{t^2},
\end{equation}
where
\begin{equation}\label{eq:kappa}
 \kappa=-\frac6{\pi^2}\int_0^\infty
 \frac{y e^{-3y/2}}{1-e^{-y}}
 \int_0^1\ell(x)
 \left[\cosh\!\left(y\left(x-\frac12\right)\right)-1\right]
 \dd x\dd y>0.
\end{equation}
\end{proposition}

\begin{proof}
Insert \eqref{eq:Delta} into \eqref{eq:mathcalD}.  With
\[
 F(y,x)=\frac{e^{-3y/2}}{1-e^{-y}}\ell(x)
 \bigl(e^{-y(x-1/2)}-1\bigr),
\]
one obtains
\[
 t^2\mathcal D(t)
 =t^2\sum_{k\geq2}\sum_{\substack{1\leq r<k\\(r,k)=1}}
 F(tk,r/k).
\]
For fixed \(0<\delta<1/2\) and \(0<\eta<M<\infty\), truncate to
\(\delta\leq r/k\leq1-\delta\) and \(\eta\leq tk\leq M\).  On this
compact region the weight is bounded and continuous, so the classical
primitive-lattice Riemann-sum limit obtained by M\"obius inversion gives
\[
 t^2\sum_{k\geq1}\sum_{\substack{1\leq r<k\\(r,k)=1}}
 F(tk,r/k)
 \longrightarrow\frac6{\pi^2}\int_0^\infty y
 \int_0^1F(y,x)\dd x\dd y
\]
after the three truncations are removed; see
\cite[Chapter~XVIII]{HardyWright2008}.  We now justify that removal.

For \(0<x\leq1/2\),
\begin{equation}\label{eq:ell-bound-D}
 |\ell(x)|\ll1+\log(1/x),
\end{equation}
and the same bound holds near \(1\) by reflection.  Also
\[
 \left|\frac{e^{-3y/2}}{1-e^{-y}}
 \bigl(e^{-y(x-1/2)}-1\bigr)\right|
 \leq \frac{y e^{-y}}{2(1-e^{-y})}=:h(y),
\]
where \(h(y)=O(1)\) as \(y\downarrow0\) and
\(h(y)\ll y e^{-y}\) as \(y\to\infty\).  From
\eqref{eq:ell-bound-D} and Stirling's estimate for
\(\sum_{r\leq \delta k}\log r\), uniformly in \(k\),
\begin{equation}\label{eq:boundary-discrete}
 \sum_{1\leq r\leq\delta k}|\ell(r/k)|
 \ll k\delta\bigl(1+|\log\delta|\bigr)+\log(2k).
\end{equation}
Hence the two endpoint strips contribute, in the limit superior as
\(t\downarrow0\), at most
\(O(\delta(1+|\log\delta|))\); the term containing \(\log(2k)\)
is \(o(1)\).

For the \(y\)-tails, the complete sine product gives the exact bound
\[
 \sum_{r=1}^{k-1}|\ell(r/k)|
 =(k-1)\log2-\log k\ll k.
\]
Therefore the range \(tk<\eta\) contributes \(O(\eta^2)\), while the
range \(tk>M\) contributes
\[
 O\!\left(\int_M^\infty y h(y)\dd y\right)
 =O\!\left((1+M)^2e^{-M}\right),
\]
uniformly for sufficiently small \(t\).  Let first \(t\downarrow0\), then
\(\delta,\eta\downarrow0\), and finally \(M\to\infty\).  This proves the
untruncated limit despite the logarithmic endpoint singularities.

Finally, the symmetry \(\ell(x)=\ell(1-x)\) replaces the exponential
bracket by the hyperbolic cosine in \eqref{eq:kappa}.  The resulting
integral is absolutely convergent by the same bounds.  Since
\(\ell(x)<0\) for \(0<x<1\), except at \(x=1/2\), and
\(\cosh(y(x-1/2))-1\geq0\), the inner integral in
\eqref{eq:kappa} is strictly negative for every \(y>0\); hence
\(\kappa>0\).
\end{proof}

\section{Incomplete reduced-residue sums}
\label{sec:clausen}

For \(0\leq S<m\), put
\begin{equation}\label{eq:Hpartial}
 \mathcal H(m,S)=\sum_{s=1}^{S}
 \log\!\left(2\sin\frac{\pi s}{m}\right),
 \qquad
 \operatorname{Cl}_2(\theta)=\sum_{j\geq1}\frac{\sin(j\theta)}{j^2}.
\end{equation}
The Clausen function and its relation to the dilogarithm are classical;
see \cite{Lewin1981}.

\begin{proposition}\label{prop:moving-sieve}
For \(q>1\) and \(0\leq R<q\),
\begin{equation}\label{eq:moving-sieve}
 \sum_{\substack{1\leq r\leq R\\(r,q)=1}}
 \log\!\left(2\sin\frac{\pi r}{q}\right)
 =\sum_{e\mid q}\mu(e)\,
 \mathcal H\!\left(\frac qe,\left\lfloor\frac Re\right\rfloor\right).
\end{equation}
\end{proposition}

\begin{proof}
Insert \(\mathbf1_{(r,q)=1}=\sum_{e\mid(r,q)}\mu(e)\) and write
\(r=es\).  Then \(r\leq R\) is equivalent to
\(s\leq\lfloor R/e\rfloor\), while \(r/q=s/(q/e)\).
\end{proof}

The floor in \eqref{eq:moving-sieve} is essential:
\begin{equation}\label{eq:floor-phase}
 \frac{\lfloor R/e\rfloor}{q/e}
 =\frac{R-(R\bmod e)}q.
\end{equation}
Thus the Clausen phases are not exactly independent of \(e\).

\begin{lemma}\label{lem:partial-sine}
For \(1\leq S<m\), with \(a=S/m\),
\begin{align}
 \mathcal H(m,S)
 &=-\frac{m}{2\pi}\operatorname{Cl}_2(2\pi a)
 +\frac12\log(2\pi S)
 +\frac12\log\!\left(\frac{\sin(\pi a)}{\pi a}\right)
 +E(m,S),\label{eq:partial-sine}\\
 |E(m,S)|
 &\leq\frac1{12S}
 +\frac12\log\!\left(\frac{\pi a}{\sin(\pi a)}\right).
 \label{eq:partial-sine-error}
\end{align}
\end{lemma}

\begin{proof}
Split
\[
 \log(2\sin\pi x)=\log(2\pi x)
 +g(x),\qquad g(x)=\log\!\left(\frac{\sin\pi x}{\pi x}\right).
\]
Stirling's formula with its elementary remainder treats the first term.
For the second, first-order Euler summation gives
\[
 \sum_{s=1}^{S}g(s/m)
 =m\int_0^ag(x)\dd x+\frac12g(a)+E_g,
 \qquad |E_g|\leq\frac12\int_0^a|g'(x)|\dd x=-\frac12g(a).
\]
Finally,
\[
 \int_0^S\log\!\left(2\sin\frac{\pi t}{m}\right)\dd t
 =-\frac{m}{2\pi}\operatorname{Cl}_2(2\pi S/m).
\]
These are the standard Euler--Maclaurin ingredients for partial rational
sine products; compare \cite{OSullivan2016}.
\end{proof}

\begin{theorem}\label{thm:Clausen}
For \(q>1\) and \(0\leq R<q\),
\begin{equation}\label{eq:Clausen}
 \sum_{\substack{1\leq r\leq R\\(r,q)=1}}
 \log\!\left(2\sin\frac{\pi r}{q}\right)
 =-\frac{\varphi(q)}{2\pi}
 \operatorname{Cl}_2\!\left(2\pi\frac Rq\right)
 +\frac12\Lambda(q)+\eps(q,R),
\end{equation}
where, uniformly in \(R\),
\begin{equation}\label{eq:epsilon-bound}
 |\eps(q,R)|\ll d(q)\log(2q).
\end{equation}
Moreover,
\begin{equation}\label{eq:epsilon-mean}
 \sum_{R=0}^{q-1}\eps(q,R)=0.
\end{equation}
\end{theorem}

\begin{proof}
For \(R=0\), take \(\eps(q,0)=-\Lambda(q)/2\).  Suppose \(R\geq1\).
Apply Lemma~\ref{lem:partial-sine} to every divisor \(e\mid q\) with
\(e\leq R\), and put
\(a_e=e\lfloor R/e\rfloor/q\) and \(a=R/q\).  Since
\(|a-a_e|<e/q\), the classical modulus estimate
\[
 |\operatorname{Cl}_2(2\pi x)-\operatorname{Cl}_2(2\pi y)|
 \ll |x-y|\log\frac{2}{|x-y|}
\]
follows by integrating
\(\operatorname{Cl}_2'(\theta)=-\log|2\sin(\theta/2)|\).  It shows that
replacing \(a_e\) by \(a\), after multiplication by \(q/e\),
costs \(O(1+\log(q/e))\) per divisor.  The divisors \(e>R\) satisfy the
same total bound because
\(|\operatorname{Cl}_2(2\pi a)|\ll a\log(2/a)\).  The Clausen terms then
collapse by
\(\sum_{e\mid q}\mu(e)/e=\varphi(q)/q\).

For \(e\leq R\),
\(\log\lfloor R/e\rfloor=\log R-\log e+O(1)\).  Completing the divisor
sums and using the known identities
\[
 \sum_{e\mid q}\mu(e)=0,
 \qquad \sum_{e\mid q}\mu(e)\log e=-\Lambda(q)
\]
extracts \(\Lambda(q)/2\).  The sinc terms and the remainders in
\eqref{eq:partial-sine-error} contribute at most
\(O(d(q)\log(2q))\).  This proves \eqref{eq:epsilon-bound}.

For the mean, pair \(r\) with \(q-r\) and use the complete cyclotomic
product to obtain
\[
 \sum_{R=0}^{q-1}\sum_{\substack{r\leq R\\(r,q)=1}}
 \log\!\left(2\sin\frac{\pi r}{q}\right)=\frac q2\Lambda(q).
\]
Also \(\sum_{R=0}^{q-1}\operatorname{Cl}_2(2\pi R/q)=0\).  Summing
\eqref{eq:Clausen} over \(R\) gives \eqref{eq:epsilon-mean}.
\end{proof}

At the moving cutoffs relevant to the diagonal product, define
\begin{equation}\label{eq:mathcalE}
 \mathcal E(X)=\sum_{q=2}^{X-1}\eps\bigl(q,X\bmod q\bigr).
\end{equation}
The uniform bound \eqref{eq:epsilon-bound} and the elementary estimate
\(\sum_{q\leq X}d(q)\ll X\log X\) give only
\begin{equation}\label{eq:mathcalE-bound}
 \mathcal E(X)\ll X(\log X)^2.
\end{equation}
Thus the pointwise Clausen decomposition does not, by direct summation of
its remainders, reach the \(\sqrt X\)-scale; any sharper use of the formula
must exploit cancellation or additional structure in the moving
remainders.

The totient-weighted Clausen aggregate has an elementary quadratic mean.

\begin{proposition}\label{prop:Clausen-main}
As \(X\to\infty\) through the integers,
\begin{equation}\label{eq:Clausen-main}
 \sum_{q<X}\varphi(q)
 \operatorname{Cl}_2\!\left(2\pi\left\{\frac Xq\right\}\right)
 =c_{\rm Cl}X^2+o(X^2),
\end{equation}
where
\begin{equation}\label{eq:cCl}
 c_{\rm Cl}=\frac6{\pi^2}\int_0^1
 \operatorname{Cl}_2(2\pi t)
 \sum_{m\geq1}\frac{\dd t}{(m+t)^3}.
\end{equation}
\end{proposition}

\begin{proof}
Set \(f(t)=\operatorname{Cl}_2(2\pi t)\) and
\(g(u)=f(\{1/u\})\).  Since \(f(0)=f(1)=0\), the bounded function \(g\)
is Riemann integrable on \([0,1]\).  Partial summation with the classical
estimate
\(\sum_{n\leq y}\varphi(n)=3y^2/\pi^2+O(y\log y)\) gives
\[
 \frac1{X^2}\sum_{q<X}\varphi(q)g(q/X)
 \longrightarrow\frac6{\pi^2}\int_0^1ug(u)\dd u.
\]
Substitute \(v=1/u\) and split \(v\) into the intervals
\([m,m+1]\) to obtain \eqref{eq:cCl}.  The mean value of \(\varphi\) is
standard; see \cite[Chapter~3]{Apostol1976}.
\end{proof}

\section{Arithmetic encodings}
\label{sec:recoveries}

This section records finite codings obtained from the zeros and complete
reduced-residue values of the product.  The prime indicator, the Boolean
M\"obius formula, and the cyclotomic recovery of \(\Lambda\) are exact
encodings; they are not analytic estimates.  The exact Chebyshev reduction
and the totient obstruction, which make genuinely different statements, are
separated into the next section.

\subsection{Prime counting and a normalized count}

Adopt the empty-product convention \(P(2)=1\).  Since \(P(n)=0\) exactly
when \(n\) has a divisor \(k\) with \(2\leq k<n\),
\begin{equation}\label{eq:prime-count}
 \pi(X)=1+\sum_{n=3}^{X}\left\lceil|P(n)|\right\rceil.
\end{equation}
This is an indicator formula of the same general kind as other exact but
computationally inefficient prime formulas; compare Willans
\cite{Willans1964}.

The exponential scale supplies a numerically stable normalization.  Put
\begin{equation}\label{eq:critical-beta}
 \beta_c=\frac{\log2}{C},
 \qquad
 S_c(X)=\sum_{2\leq n\leq X}|P(n)|^{\beta_c/n},
\end{equation}
where \(0^a=0\) for \(a>0\).  For a prime \(p\), write
\(R(p)=\log|P(p)|+Cp\).  Theorem~\ref{thm:prime-third-scale} gives
\(R(p)=O(p^{1/3}(\log p)^2)\), hence
\begin{equation}\label{eq:critical-weight}
 |P(p)|^{\beta_c/p}
 =\frac12+O\!\left(p^{-2/3}(\log p)^2\right).
\end{equation}
The elementary Chebyshev bound \(\pi(x)\ll x/\log x\) and partial
summation imply
\begin{equation}\label{eq:normalized-count}
 S_c(X)=\frac12\pi(X)+O(X^{1/3}\log X).
\end{equation}
Consequently, the classical von Koch criterion may still be written as
\begin{equation}\label{eq:RH-normalized-count}
 \mathrm{RH}\quad\Longleftrightarrow\quad
 2S_c(X)-\Li(X)=O(\sqrt X\log X).
\end{equation}
Indeed,
\begin{equation}\label{eq:RH-normalized-difference}
 2S_c(X)-\Li(X)
 =\pi(X)-\Li(X)+O(X^{1/3}\log X),
\end{equation}
and the added term is \(o(\sqrt X\log X)\).  Thus the strengthened
product remainder sharpens the approximation of \(S_c\) to \(\pi/2\), but
it does not sharpen or alter the von Koch criterion itself; see
\cite[Chapter~XIV]{TitchmarshHeathBrown1986}.

\subsection{A Boolean formula for the M\"obius function}

Define the Boolean quantities
\begin{equation}\label{eq:ABC}
 A(k)=\lceil|P(k)|\rceil,
 \quad
 B_k(n)=\left\lceil\left|\sin\frac{\pi n}{k}\right|\right\rceil,
 \quad
 C_k(n)=\left\lceil\left|\sin\frac{\pi n}{k^2}\right|\right\rceil.
\end{equation}

\begin{proposition}\label{prop:mobius-boolean}
For every \(n\geq1\),
\begin{equation}\label{eq:mobius-boolean}
 \mu(n)=\prod_{k=2}^{n}
 \left[1-A(k)+A(k)C_k(n)\bigl(2B_k(n)-1\bigr)\right].
\end{equation}
\end{proposition}

\begin{proof}
If \(k\) is composite, then \(A(k)=0\) and the \(k\)-factor is \(1\).
If \(k\) is prime, the factor is \(1\), \(-1\), or \(0\) according as
\(k\nmid n\), \(k\mid n\) but \(k^2\nmid n\), or \(k^2\mid n\).
The product is therefore zero on nonsquarefree integers and otherwise is
\((-1)^{\omega(n)}\), the standard definition of \(\mu(n)\).
\end{proof}

Formula \eqref{eq:mobius-boolean} is a finite Boolean encoding of prime
exponents.  It does not avoid the primality and divisibility tests encoded
by its sine zeros.

\subsection{Cyclotomic recovery of the von Mangoldt function}

For \(q\geq2\), define the complete reduced-residue sum
\begin{equation}\label{eq:J}
 J(q)=\sum_{\substack{1\leq r<q\\(r,q)=1}}
 \log\sin\!\left(\frac{\pi r}{q}\right).
\end{equation}
The classical cyclotomic evaluation
\begin{equation}\label{eq:cyclotomic-Lambda}
 \prod_{\substack{1\leq r<q\\(r,q)=1}}
 2\sin\!\left(\frac{\pi r}{q}\right)
 =\Phi_q(1)=e^{\Lambda(q)}
\end{equation}
gives the known identity
\begin{equation}\label{eq:J-Lambda}
 J(q)=\Lambda(q)-\varphi(q)\log2.
\end{equation}
See \cite[Lemma~2.2]{HerreraMoree2018} and
\cite[Chapters~1--2]{Washington1997}.  Thus the complete value
\(H_q(1)\) in \eqref{eq:H-classical} recovers \(\Lambda(q)\); this is a
cyclotomic identity, not a new formula for the distribution of prime
powers.

\section{Exact Chebyshev reduction and totient obstruction}
\label{sec:chebyshev-obstruction}

The following results are not merely alternative codings of arithmetic
functions.  The first removes an auxiliary normalization exactly; the second
shows that a natural termwise asymptotic decomposition is impossible at the
scale relevant to the prime error term.

\subsection{Closed totient block and four-term identity}

Write \(\Psi(X)=\sum_{m\leq X}\Lambda(m)\).  For an integer \(X\geq2\),
let
\begin{align}
 T(X)&=\sum_{n\leq X}\varphi(n),\label{eq:T-def}\\
 \mathcal B_X&=\sum_{q=2}^{X-1}
 \left(\left\lfloor\frac Xq\right\rfloor-1\right)\varphi(q),
 &Q_X&=\sum_{q=2}^{X-1}
 \sum_{\substack{1\leq r\leq X\bmod q\\(r,q)=1}}
 \log\sin\!\left(\frac{\pi r}{q}\right).
 \label{eq:BQ}
\end{align}

\begin{proposition}[Closed form for the totient block]\label{prop:B-closed}
For every integer \(X\geq2\),
\begin{equation}\label{eq:B-closed}
 \mathcal B_X=\frac{X^2-X}{2}+1-T(X).
\end{equation}
\end{proposition}

\begin{proof}
The classical divisor identity \(\sum_{d\mid m}\varphi(d)=m\) gives
\[
 \sum_{q\leq X}\left\lfloor\frac Xq\right\rfloor\varphi(q)
 =\sum_{m\leq X}\sum_{q\mid m}\varphi(q)
 =\sum_{m\leq X}m=\frac{X(X+1)}2.
\]
Removing the terms \(q=1\) and \(q=X\), and then subtracting
\(\sum_{q=2}^{X-1}\varphi(q)=T(X)-1-\varphi(X)\), yields
\eqref{eq:B-closed} exactly.
\end{proof}

\begin{proposition}\label{prop:Chebyshev}
For every \(X\geq2\),
\begin{equation}\label{eq:Chebyshev}
 \boxed{\displaystyle
 \Psi(X)=\log(X!)-(\log2)\,\mathcal B_X+Q_X
 -\sum_{n\leq X}L(n).}
\end{equation}
In particular, the normalization constant \(C\) cancels identically.
\end{proposition}

\begin{proof}
Sum \(L(n)\) first over \(n\leq X\) and then over its denominator \(q\).
For fixed \(q\), the residues consist of
\(\lfloor X/q\rfloor-1\) complete reduced-residue cycles and a terminal
segment.  Formula \eqref{eq:J-Lambda} gives
\[
 \sum_{n\leq X}L(n)
 =\sum_{q=2}^{X-1}
 \left(\left\lfloor\frac Xq\right\rfloor-1\right)
 \bigl(\Lambda(q)-\varphi(q)\log2\bigr)+Q_X.
\]
The standard identity
\[
 \sum_{q\leq X}\left\lfloor\frac Xq\right\rfloor\Lambda(q)
 =\log(X!)
\]
follows by summing \(\log m=\sum_{q\mid m}\Lambda(q)\) over
\(m\leq X\).  Hence
\[
 \sum_{q=2}^{X-1}
 \left(\left\lfloor\frac Xq\right\rfloor-1\right)\Lambda(q)
 =\log(X!)-\Psi(X),
\]
and rearrangement gives \eqref{eq:Chebyshev}.  Equivalently, in the
normalization \(R_L=L+C\varphi\), one has
\(\sum_{n\leq X}R_L(n)=\sum_{n\leq X}L(n)+CT(X)\), so the two
\(CT(X)\) terms in the previous normalized form cancel exactly.
\end{proof}

Combining Propositions~\ref{prop:B-closed} and \ref{prop:Chebyshev} gives
the four-term form
\begin{equation}\label{eq:Chebyshev-four-term}
 \Psi(X)=\log(X!)-(\log2)\frac{X^2-X+2}{2}
 +\left[(\log2)T(X)+Q_X-\sum_{n\leq X}L(n)\right].
\end{equation}
This form makes an unavoidable cancellation visible.

\subsection{A proved obstruction to termwise estimation}

\begin{proposition}[Totient obstruction to termwise estimation]
\label{prop:totient-obstruction}
Let
\begin{equation}\label{eq:Ephi}
 E_\varphi(X)=T(X)-\frac3{\pi^2}X^2.
\end{equation}
Then
\begin{equation}\label{eq:Ephi-Omega}
 \limsup_{X\to\infty}\frac{|E_\varphi(X)|}{X}>0.
\end{equation}
Consequently, if
\begin{equation}\label{eq:MQ}
 M_Q(X)=X-\log(X!)+(\log2)
 \left(\frac{X^2-X+2}{2}-\frac3{\pi^2}X^2\right),
\end{equation}
then
\begin{equation}\label{eq:QminusL-Omega}
 \limsup_{X\to\infty}\frac1X
 \left|Q_X-\sum_{n\leq X}L(n)-M_Q(X)\right|>0.
\end{equation}
Thus the \(\Omega(X)\) fluctuation in
\(Q_X-\sum_{n\leq X}L(n)\) is forced by the totient term; it is larger
than the \(O(\sqrt X\log^2X)\) precision required by the von Koch form of
the Riemann-hypothesis error term for \(\Psi\).
\end{proposition}

\begin{proof}
For a prime \(p\),
\begin{align*}
 E_\varphi(p)-E_\varphi(p-1)
 &=\varphi(p)-\frac3{\pi^2}\bigl(p^2-(p-1)^2\bigr)\\
 &=\left(1-\frac6{\pi^2}\right)p+O(1).
\end{align*}
The coefficient is positive.  Therefore at least one of
\(|E_\varphi(p)|\) and \(|E_\varphi(p-1)|\) is at least a fixed positive
multiple of \(p\) for every sufficiently large prime \(p\), proving
\eqref{eq:Ephi-Omega}.

From \eqref{eq:Chebyshev-four-term},
\begin{equation}\label{eq:QminusL-decomp}
 Q_X-\sum_{n\leq X}L(n)-M_Q(X)
 =\Psi(X)-X-(\log2)E_\varphi(X).
\end{equation}
The prime number theorem gives \(\Psi(X)-X=o(X)\), whereas
\eqref{eq:Ephi-Omega} supplies an \(\Omega(X)\) term.  This proves
\eqref{eq:QminusL-Omega}.  In particular, replacing \(T(X)\) by its
quadratic main term before combining it with
\(Q_X-\sum_{n\leq X}L(n)\) necessarily leaves an error of order \(X\)
along a subsequence.  The cancellation in the bracket of
\eqref{eq:Chebyshev-four-term} must therefore be retained.
\end{proof}

Thus \eqref{eq:Chebyshev} remains a finite reorganization of the classical
divisor identity for \(\Lambda\), but \eqref{eq:QminusL-Omega} identifies
precisely where a term-by-term asymptotic treatment fails: the dominant
uncancelled fluctuation is totient rather than prime.

\section{Numerical audit}
\label{sec:numerical-audit}

\begin{experiment}\label{exp:audit}
All computations in this section check finite identities, constants, and
finite-range conjectural diagnostics.  None is used in a proof.  The
standalone script \texttt{v8\_numeric\_audit.py} reproduces the new tables.
\end{experiment}

\subsection{Previously established finite checks}

Independent quadrature of \eqref{eq:C-integrals} and summation of
\eqref{eq:C-series} agree to \(7\times10^{-13}\), with
\[
 C=0.7437948618394\ldots,
 \qquad C-\log2=0.0506476813\ldots.
\]
The divisor decomposition \eqref{eq:primitive} was checked for every
\(n\le1200\); the largest floating-point discrepancy was
\(2.3\times10^{-13}\).

For the exact cumulative-frequency object \eqref{eq:BJ-quadrature}, direct
evaluation at \(p=100003\) gives
\[
 \frac{B_{1000}}{\sqrt{1000p}}=-0.269309229,
 \quad
 \frac{B_{2000}}{\sqrt{2000p}}=-0.578604231,
 \quad
 \frac{B_{3000}}{\sqrt{3000p}}=-0.427901820.
\]
These values are retained only as a diagnostic of the exact quadrature
identity; Section~\ref{sec:frequency-centered} no longer uses a square-root
\(B_J\) criterion as the best pointwise route.

The two sides of the Lambert identity \eqref{eq:Lambert}, truncated at
index \(200\), agreed to \(8.7\times10^{-19}\) at \(z=0.3\) and to
\(6.7\times10^{-16}\) at \(z=0.7\).  Direct quadrature gives
\(\kappa=0.05826334\ldots\).  The exact sieve
\eqref{eq:moving-sieve} was checked for all \(2\le q\le300\) and
\(0\le R<q\), with maximum discrepancy \(2.18\times10^{-14}\).  The
Boolean identity \eqref{eq:mobius-boolean} and prime count
\eqref{eq:prime-count} were checked through \(n=20000\); the closed totient
form \eqref{eq:B-closed} has zero integer discrepancy for every
\(3\le X\le300\), and the simplified Chebyshev identity was checked at
\(X=50,100,200,300\) to \(1.7\times10^{-12}\).

\subsection{Smooth/integer split and gamma audit}

Quadrature of \eqref{eq:J0-def}--\eqref{eq:H0-def} gives
\[
 J_0=-1.750632969245044\ldots,
 \qquad H_0=-0.137891778443756\ldots,
\]
and
\[
 \log\pi+J_0+H_0=-0.743794861839401\ldots=-C
\]
to the displayed precision.  Identity \eqref{eq:integer-split} was checked
directly at \(p=1009\) and \(p=10007\), with discrepancies
\(-4.55\times10^{-13}\) and \(5.46\times10^{-12}\), respectively.

For \(p=10007\), Proposition~\ref{prop:NpS-exact} gives exact integer
agreement at each tested cutoff:
\[
\begin{array}{c|rrrrr}
 S&3&7&15&30&60\\ \hline
 \mathcal N_p(S)&75&119&206&374&658\\
 \text{right side of \eqref{eq:NpS-exact}}&75&119&206&374&658
\end{array}
\]
The corrected gamma factorization is sensitive to the direction of the
Pochhammer quotients.  At the same prime,
\[
 \log S_p\ \text{(direct)}=64659.426286523681\ldots,
\]
while \eqref{eq:Sp-gamma} gives
\[
 64659.426286523696\ldots,
\]
a discrepancy of \(-1.46\times10^{-11}\).  The reflected factor is
therefore \(|\sin(\pi p/q)|/\pi\), not its reciprocal.

\subsection{Exact regular-grid identity for the folded discrepancy}

Equation \eqref{eq:Hp-grid} was compared with a direct strict count, not
with a reconstruction from the same formula.  Representative values are
\[
\begin{array}{c|c|r|r|r}
 p&t&H_p^{<}(t)&\text{right side of \eqref{eq:Hp-grid}}&\text{difference}\\ \hline
1009&0.073& 4.573630754732& 4.573630754732&-3.7\cdot10^{-13}\\
1009&0.123& 2.586401744951& 2.586401744951& 3.0\cdot10^{-13}\\
10007&0.073&4.575725433698&4.575725433686&1.1\cdot10^{-11}\\
10007&0.271&-3.885421385915&-3.885421385923&8.4\cdot10^{-12}\\
100003&0.123&-15.233017145323&-15.233017145309&-1.4\cdot10^{-11}\\
100003&0.271&-6.681468857360&-6.681468857427&6.7\cdot10^{-11}
\end{array}
\]
At an exact jump, a computation using \(\le\) instead of \(<\) differs by
the corresponding integer jump, as stated in
Theorem~\ref{thm:Hp-grid}.  Omitting the digamma and boundary corrections
leaves a bounded residual approaching \(2t-1\), in agreement with
\eqref{eq:Hp-grid-smooth-asymp}.

\subsection{A failed pure-area localized count}

The stronger heuristic
\[
 \#\{k\le\sqrt p:\|p/k\|<\delta\}
 \stackrel{?}{=}O(\delta\sqrt p)+O(1)
\]
is false as a natural uniform model because divisor-rich shifts create
spikes.  For
\[
 p=1441439,\qquad K=1200,\qquad \delta=2/1200,
\]
the direct count is \(30\), while \(\delta\sqrt p=2.000999\ldots\).
Here \(p+1=1441440\) has many divisors in the relevant range.  No argument
in this version uses the rejected pure-area estimate.

\subsection{Exhaustive folded-discrepancy maxima}

By Lemma~\ref{lem:Hp-jumps}, evaluating both one-sided values at every jump
is exhaustive for the continuous supremum.  On the mesoscopic interval
\(p^{-2/3}\le t\le1/2\) we obtain
\[
\begin{array}{r|r|r|r}
 p&\max|H_p|&\max|H_p|/p^{1/3}&\max|H_p|/(p^{1/3}\log p)\\ \hline
1009&7.127365&0.710611&0.102738\\
5003&15.867216&0.927735&0.108917\\
10007&14.505816&0.673143&0.073080\\
20011&19.008938&0.700167&0.070695\\
50021&24.758911&0.671966&0.062103\\
100003&40.335013&0.868983&0.075479\\
300007&41.881577&0.625623&0.049607
\end{array}
\]
The first normalized column stays between \(0.62\) and \(0.93\) on this
sample, whereas division by the additional \(\log p\) decreases overall.
This is finite evidence for \eqref{eq:Hp-conjecture}, not a proof.

At \(p=1000003\), the exact one-sided jump maxima in dyadic bands are
\[
\begin{array}{c|r@{\qquad}c|r}
[t_0,t_1)&\max|H_p|&[t_0,t_1)&\max|H_p|\\ \hline
{}[1/4,1/2)&44.5111 &[1/64,1/32)&33.8653\\
{}[1/8,1/4)&40.9914 &[1/128,1/64)&31.6146\\
{}[1/16,1/8)&41.5335 &[1/256,1/128)&31.2187\\
{}[1/32,1/16)&44.9391 &[1/512,1/256)&21.1153\\
{}[1/1024,1/512)&22.6831 &[1/2048,1/1024)&14.0273\\
{}[1/4096,1/2048)&12.3144 &[1/8192,1/4096)&7.3583
\end{array}
\]
The maximum over these displayed bands is \(44.9391\ldots\), while
\(p^{1/3}=100.0001\ldots\).

\subsection{Normalized prime-count scale}

The deterministic majorant entering \eqref{eq:normalized-count} is
\[
 \sum_{p\le X}p^{-2/3}(\log p)^2.
\]
Its values are
\[
\begin{array}{r|r|r|r}
X&\text{sum}&X^{1/3}\log X&\text{ratio}\\ \hline
10^4&400.225&198.431&2.017\\
10^5&1181.741&534.383&2.211\\
10^6&3238.226&1381.551&2.344
\end{array}
\]
This only checks the expected \(X^{1/3}\log X\) order of the auxiliary
normalization error.  It does not test the von Koch remainder
\(\pi(X)-\Li(X)\).

\section{Conclusion}
\label{sec:conclusion}

The diagonal reciprocal sine product has two quantitatively distinct
levels.  The elementary block decomposition gives the density
\(\rho(t)=\psi_1(1+t)\) and an \(O(\sqrt n\log n)\) logarithmic remainder.
The sharper Saffari--Vaughan discrepancy replaces this, on primes, by the
proved bound
\[
 \log|P(p)|=-Cp+O(p^{1/3}(\log p)^2),
\]
and propagates to the zero-omitted product, the primitive component
\(R_L(n)=O_\eps(n^{1/3+\eps})\), and absolute convergence of the associated
Dirichlet series for \(\Re s>4/3\).  The exact rotated Parseval identity
simultaneously improves to
\(\|\mathscr R_p\|_2\ll p^{1/3}\log p\).

The product-specific information begins after this imported discrepancy
estimate.  Reflection of the logarithmic sine kernel folds the full
reciprocal discrepancy into \(H_p(t)=D_p(t)-D_p(1-t)\).  The exact
Stieltjes formula \eqref{eq:folded-Stieltjes} shows that the extreme
singular range contributes only \(O(p^{1/3}\log p)\); the second logarithm
in the current pointwise theorem can arise only in the mesoscopic range
\(p^{-2/3}\le t\le1/2\).  The exact sawtooth identity
\eqref{eq:Hp-grid} puts this folded object on a regular reciprocal grid,
including the digamma and boundary corrections that are numerically
essential.  Exhaustive finite computations support, but do not prove, the
stronger bound \(H_p(t)\ll p^{1/3}\), which would remove one logarithm from
the product remainder.

The singular logarithm also separates into a bounded-variation term and the
integer core \(S_p=\prod s_k\).  The latter is an exact statistic of the
symmetric shifted divisors of \(p-s\) and \(p+s\), admits the corrected
gamma factorization \eqref{eq:Sp-gamma}, and has a low-complementary-divisor
sector that can be smoothed to precision \(O(p^{1/3}\log p)\).  Gamma
reflection alone does not close the remaining pointwise problem.

The later Euclidean Lambert, Clausen, cyclotomic, M\"obius, totient, and
Chebyshev identities remain exact structural reorganizations.  The improved
prime normalization now satisfies
\[
 S_c(X)=\frac12\pi(X)+O(X^{1/3}\log X),
\]
but this sharper auxiliary approximation is still absorbed inside the
classical von Koch scale and therefore does not yield new information about
the zeros of \(\zeta\).  The unresolved mathematical problem specific to
the present product is the mesoscopic folded discrepancy, not the Riemann
hypothesis.

\begin{thebibliography}{99}
\raggedright

\bibitem{Apostol1976}
T.~M. Apostol,
\emph{Introduction to Analytic Number Theory},
Springer, New York, 1976.
\href{https://doi.org/10.1007/978-1-4757-5579-4}
{doi:10.1007/978-1-4757-5579-4}.

\bibitem{GrahamKolesnik1991}
S.~W. Graham and G.~Kolesnik,
\emph{Van der Corput's Method of Exponential Sums},
London Mathematical Society Lecture Note Series 126,
Cambridge University Press, Cambridge, 1991.
\href{https://doi.org/10.1017/CBO9780511661976}
{doi:10.1017/CBO9780511661976}.

\bibitem{GrepstadEtAl2019}
S.~Grepstad, L.~Kaltenb\"ock, and M.~Neum\"uller,
\emph{On the asymptotic behaviour of the sine product
\(\prod_{r=1}^n|2\sin(\pi r\alpha)|\)},
arXiv:1909.00980 (2019).
\url{https://arxiv.org/abs/1909.00980}.

\bibitem{HardyLittlewoodPolya1952}
G.~H. Hardy, J.~E. Littlewood, and G.~P\'olya,
\emph{Inequalities}, 2nd ed., Cambridge University Press, Cambridge, 1952.

\bibitem{HardyWright2008}
G.~H. Hardy and E.~M. Wright,
\emph{An Introduction to the Theory of Numbers}, 6th ed.,
Oxford University Press, Oxford, 2008.

\bibitem{HerreraMoree2018}
J.~Herrera-Poyatos and P.~Moree,
\emph{Coefficients and higher order derivatives of cyclotomic polynomials:
old and new}, MPIM Preprint 2018 (46).
\url{https://archive.mpim-bonn.mpg.de/1354/}.

\bibitem{Lewin1981}
L.~Lewin,
\emph{Polylogarithms and Associated Functions},
North-Holland, New York, 1981.

\bibitem{LiYang2023}
X.~Li and X.~Yang,
\emph{An improvement on Gauss's circle problem and Dirichlet's divisor problem},
arXiv:2308.14859 (2023), preprint.
\url{https://arxiv.org/abs/2308.14859}.

\bibitem{Mayer1989}
D.~H. Mayer,
\emph{Continued fractions and related transformations},
MPIM Preprint 1989 (73).
\url{https://archive.mpim-bonn.mpg.de/1282/}.

\bibitem{OSullivan2016}
C.~O'Sullivan,
\emph{Asymptotics for the partial fractions of the restricted partition
generating function I},
International Journal of Number Theory \textbf{12} (2016), 1421--1474.
\href{https://doi.org/10.1142/S1793042116500895}
{doi:10.1142/S1793042116500895}.

\bibitem{MikeTeX2023}
MikeTeX,
\emph{Distribution of the fractional parts of \(n/1,n/2,\ldots,n/n\) as
\(n\) tends to infinity},
Mathematics Stack Exchange, question 4831482, 21 December 2023.
\url{https://math.stackexchange.com/questions/4831482/}.

\bibitem{SaffariVaughan1977}
B.~Saffari and R.~C. Vaughan,
\emph{On the fractional parts of \(x/n\) and related sequences. II},
Annales de l'Institut Fourier \textbf{27} (1977), no.~2, 1--30.
\href{https://doi.org/10.5802/aif.649}{doi:10.5802/aif.649}.

\bibitem{Schoenberg1938}
I.~J. Schoenberg,
\emph{Metric spaces and positive definite functions},
Transactions of the American Mathematical Society \textbf{44} (1938),
522--536.
\href{https://doi.org/10.1090/S0002-9947-1938-1501980-0}
{doi:10.1090/S0002-9947-1938-1501980-0}.

\bibitem{TitchmarshHeathBrown1986}
E.~C. Titchmarsh,
\emph{The Theory of the Riemann Zeta-Function}, 2nd ed.,
revised by D.~R. Heath-Brown, Clarendon Press, Oxford, 1986.

\bibitem{Washington1997}
L.~C. Washington,
\emph{Introduction to Cyclotomic Fields}, 2nd ed.,
Graduate Texts in Mathematics 83, Springer, New York, 1997.
\href{https://doi.org/10.1007/978-1-4612-1934-7}
{doi:10.1007/978-1-4612-1934-7}.

\bibitem{Willans1964}
C.~P. Willans,
\emph{On formulae for the nth prime number},
The Mathematical Gazette \textbf{48} (1964), no.~366, 413--415.
\href{https://doi.org/10.2307/3611701}{doi:10.2307/3611701}.

\end{thebibliography}

\end{document}
